% ---------------------------------------------------------------------------
% Perishable inventory section.
%
% INPUT-ABLE FRAGMENT: this file is \input{} into the manuscript
% learning_inventory_control_policies_es.tex and reuses ITS preamble macros
% and TikZ styles (envbox, envsupply/envstock/envsink/envlost,
% reqflow/matflow/demandflow/lostflow, ltlabel/reqlabel, costH/costP/costK,
% the boxed "Per-period event sequence" panel, booktabs). It defines NO new
% styles and contains NO \documentclass.
%
% ALGORITHMIC / DATA-PROVENANCE NOTE (every number is sourced from the repo):
%   Instance + dynamics : src/problems/perishable_inventory/env.rs (step_state,
%       apply_demand_to_inventory FIFO/LIFO, waste = oldest leftover bucket),
%       flownet/formulation.rs (event schedule: serve demand -> expire oldest ->
%       age + inbound receipt), references.rs (de_moor2022_m2_exp2_l1_cp7_fifo
%       and de_moor2022_m2_exp1_l1_cp7_lifo: mu=4, cov=0.5 rounded Gamma,
%       shelf life m=2, lead time L=1, h=1, shortage p=5, waste cp=7,
%       procurement c=3, order cap 10, horizon 465, warm-up 100).
%   Objective sign      : rollout.rs discounted_return_after_warmup returns
%       sum_t (-cost_t) gamma^t over the post-warm-up window (gamma=0.99), so
%       reported numbers are NEGATED discounted costs (higher / less negative is
%       better) -- the De Moor (2022) / Farrington (2025) return convention.
%   Results             : outputs/autoresearch/perishable_inventory_autoresearch/
%       results.tsv and the two *_d2_oblique_full.json files (commit 01c657a;
%       committed at 222d46f). FIFO learned -1457.9032, gate -1475.0825
%       (+1.16%, paired SEM 1.815 -> 9.5x), LIFO learned -1553.1608,
%       gate -1565.9833 (+0.82%, paired SEM 1.949 -> 6.6x). VI optima
%       -1457.28 / -1552.99 (rounded -1457 / -1553) are the ANALYTIC reference,
%       context only (estimator mismatch, literature/README.md).
% ---------------------------------------------------------------------------
\section{Perishable inventory}
\label{sec:perishable}

\subsection{Problem formulation}
\label{sec:perishable-formulation}

Figure~\ref{fig:env-perishable} depicts this environment as a flow network --- the inbound order and
its lead-time shipment into the youngest age bucket, the demand draw served by the FIFO or LIFO
issuing rule, the outdating of the oldest leftover stock, and where each cost is charged.

\begin{figure}[t]
  \centering
  \resizebox{\textwidth}{!}{%
  \begin{tikzpicture}[x=1cm,y=1cm]
    % Single-location perishable system drawn as an event-based interaction. The
    % age-structured stock is one node carrying the m=2 age buckets (youngest |
    % oldest). Inbound request UP (dashed control), lead-time arrival into the
    % youngest bucket DOWN (solid material). Demand served by FIFO/LIFO; the
    % oldest leftover bucket outdates to a waste sink; remaining stock ages.
    \node[envsupply] (sup)  at (0,0)     {external\\supplier};
    \node[envstock]  (stk)  at (7.4,0)   {age-structured stock\\$\mathbf S_t=(o^{1}_t,\,o^{2}_t)$\\(youngest $\mid$ oldest)};
    \node[envsink]   (dem)  at (15.4,0)  {rounded-Gamma\\demand $D_t$};
    \node[envlost]   (lost) at (15.4,-4.2){unmet demand\\\emph{lost}};
    \node[envlost,draw=myorange!60!black,fill=myorange!8] (waste) at (7.4,-4.2)
          {oldest leftover\\\emph{outdates} (waste $w_t$)};

    % inbound: dashed request UP (event 2), solid lead-time arrival DOWN into the
    % youngest bucket (event 5, placed L periods ago; for L=1 it arrives next period).
    \draw[reqflow] (stk.north west) to[bend right=34] (sup.north east);
    \node[reqlabel] at (3.3,2.05) {\ding{183}\, order $q_t\!\le\!\bar q$};
    \draw[matflow] (sup.south east) to[bend right=10] (stk.west);
    \node[ltlabel]  at (3.5,-1.55) {\ding{187}\, arrival (placed $L$ ago)\\enters youngest bucket; cost $c\,q$};

    \draw[demandflow] (stk) -- (dem)
          node[midway,above=4pt,ltlabel] {\ding{185}\, serve\\\textbf{FIFO} oldest-first\\\textbf{LIFO} youngest-first};
    \draw[lostflow] (dem) -- (lost);
    \draw[lostflow,myorange!70!black] (stk.south) -- (waste.north);

    \node[costH,above=9mm of stk.north] {holding $h\,(o^{1}_t)$};
    \node[costP,right=3pt of lost.east] {lost-sales\\penalty $p\,\ell_t$};
    \node[costK,below=2pt of waste.south,align=center] {waste $c_w\,w_t$};

    \begin{scope}[on background layer]
      \node[envbox,fit=(current bounding box)] (box) {};
    \end{scope}
    \node[envtitle,above=3pt of box.north,anchor=south]
          {Perishable single-item environment, shelf life $m=2$ (one period)};
  \end{tikzpicture}%
  }
  \par\vspace{6pt}\centerline{%
    \tikz\node[draw=mydarkblue!55,rounded corners=4pt,fill=mydarkblue!4,inner sep=7pt]{%
      \begin{minipage}{\dimexpr\linewidth-18pt\relax}\centering
        {\small\bfseries\color{mydarkblue!85}Per-period event sequence}\\[4pt]
        {\small\begin{tabular}{@{}p{4.5cm}p{4.5cm}p{4.5cm}@{}}
          \ding{182}\,observe age-structured state $\mathbf S_t=(o^{1}_t,o^{2}_t)$ & \ding{183}\,order $q_t\in\{0,\dots,\bar q\}$ & \ding{184}\,demand $D_t$ drawn (rounded Gamma)\\[4pt]
          \ding{185}\,serve $D_t$ by FIFO (oldest first) / LIFO (youngest first); unmet \emph{lost} & \ding{186}\,oldest leftover bucket \emph{outdates} to waste; survivors age & \ding{187}\,arrival enters youngest bucket; charge $c\,q_t+h\,o^{1}_t+p\,\ell_t+c_w w_t$\\
        \end{tabular}}
      \end{minipage}};}\par
  \caption{Perishable single-item environment of \citet{deMoor2022perishable}, drawn as an event-based
  interaction (shelf life $m=2$, lead time $L=1$). Each period the stock point \emph{sends an order
  request} $q_t$ (thin dashed control arrow); the external supplier ships and the goods \emph{arrive}
  after the lead time $L$ to enter the \emph{youngest} age bucket (solid material arrow; unit cost
  $c$). Demand $D_t$ is served from the on-hand age buckets under the configurable \emph{issuing
  rule} --- FIFO issues the \emph{oldest} stock first, LIFO the \emph{youngest} first --- and any
  \emph{unmet demand is lost}. After demand, the \emph{oldest leftover bucket outdates} to a waste
  sink (orange) and the survivors age by one bucket. Holding cost $h$ is charged on the non-expiring
  on-hand stock, each lost unit a penalty $p$, each outdated unit a waste cost $c_w$, and each ordered
  unit a procurement cost $c$; the boxed panel beneath gives the per-period event sequence
  (Section~\ref{sec:perishable-formulation}, Eq.~\eqref{eq:perish-cost}).}
  \label{fig:env-perishable}
\end{figure}

We consider the single-item, periodic-review perishable inventory model of
\citet{deMoor2022perishable}, as re-derived and benchmarked by \citet{farrington2025perishable}.
A single item with a \emph{fixed shelf life} of $m$ periods is replenished from one supplier with
deterministic lead time $L\ge 1$. Demand $D_t$ is i.i.d.\ across periods, obtained by sampling a
Gamma distribution with mean $\mu$ and coefficient of variation $\mathrm{cv}$ and rounding to the
nearest non-negative integer. We study the two instances that carry an in-crate exact-MDP verifier:
both have $m=2$, $L=1$, $\mu=4$, $\mathrm{cv}=0.5$, and differ only in the issuing rule, FIFO or LIFO.
Unmet demand is \emph{lost}. We write $(x)^+=\max(0,x)$.

\paragraph{Timing of a period.}
On-hand stock is tracked by remaining shelf life. At the start of period $t$ the controller observes
the age-structured stock, places an order $q_t$, demand $D_t$ is realized and served from on-hand
stock under the issuing rule (FIFO issues oldest units first, LIFO youngest first), unmet demand is
\emph{lost}, the \emph{oldest leftover units then outdate} (expire and are scrapped), the surviving
units \emph{age} by one bucket, and the order placed $L$ periods earlier arrives into the
\emph{youngest} bucket. With $L=1$ the order $q_t$ arrives at the start of period $t+1$ and there is
no in-transit pipeline.

\paragraph{State.}
Because shelf life is fixed at $m$ buckets and only $L-1$ orders are in transit, the MDP state is the
on-hand age vector together with the in-transit pipeline,
\begin{equation*}
    \mathbf{S}_t=\bigl(\underbrace{q_{t-L+1},\dots,q_{t-1}}_{\text{pipeline (}L-1\text{ entries)}};\ \underbrace{o^{1}_t,\dots,o^{m}_t}_{\text{on-hand by age}}\bigr),
\end{equation*}
where $o^{j}_t$ is the on-hand quantity with $j$ periods elapsed ($o^{1}_t$ youngest, $o^{m}_t$
oldest). For the studied slice $m=2,\,L=1$ the pipeline is empty and the state collapses to the
two-dimensional age vector $\mathbf{S}_t=(o^{1}_t,o^{2}_t)$.

\paragraph{Action.}
The decision in period $t$ is the scalar order quantity
\begin{equation*}
    a_t=q_t\in\{0,1,\dots,\bar q\},
\end{equation*}
where $\bar q$ is the per-period order cap.

\paragraph{Transition.}
Let $O_t=\sum_{j}o^{j}_t$ be the total on-hand stock. Demand consumes $\min(D_t,O_t)$ units under the
issuing rule, leaving a lost quantity $\ell_t=(D_t-O_t)^+$. Writing $r^{j}_t$ for the leftover in
bucket $j$ after issuing, the oldest leftover outdates,
\begin{equation*}
    w_t \;=\; r^{m}_t \qquad\text{(outdated units)},
\end{equation*}
the surviving buckets age by one slot, and the arrival $q_{t-L+1}$ (for $L=1$, the order $q_t$ just
placed) enters the youngest bucket:
\begin{equation*}
    o^{1}_{t+1}=\text{arrival},\qquad o^{j+1}_{t+1}=r^{j}_t\ \ (j=1,\dots,m-1).
\end{equation*}

\paragraph{One-period cost.}
The cost incurred in period $t$ is the linear procurement cost of the order plus holding cost on the
non-expiring on-hand stock, the lost-sales penalty, and the outdating (waste) cost:
\begin{equation}
    \label{eq:perish-cost}
    c_t \;=\; c\,q_t \;+\; h\sum_{j=1}^{m-1} r^{j}_t \;+\; p\,\ell_t \;+\; c_w\,w_t ,
\end{equation}
where $c$ is the unit procurement cost, $h$ the holding cost per unit of surviving on-hand stock, $p$
the lost-sales penalty per unit short, and $c_w$ the waste cost per outdated unit. For the studied
instances $c=3$, $h=1$, $p=5$, and the waste cost is $c_w=7$.

\paragraph{Objective.}
A stationary policy $\pi_{\btheta}$ maps the state $\mathbf{S}_t$ to the order $q_t$. Following the
discounted-return convention of \citet{deMoor2022perishable} and \citet{farrington2025perishable},
performance is the expected discounted \emph{negated} cost over the post-warm-up evaluation window
(discount factor $\gamma$, warm-up of $W$ periods within a horizon of $T$),
\begin{equation}
    \label{eq:perish-objective}
    G(\btheta) \;=\; \mathbb{E}\!\left[\sum_{t=W+1}^{T}\gamma^{\,t-W-1}\bigl(-c_t\bigr)\right],
    \qquad \max_{\btheta}\ G(\btheta),
\end{equation}
so that reported returns are \emph{negative} (lower cost is a larger, less-negative return) and higher
is better. We use $\gamma=0.99$, $T=465$, and $W=100$, i.e.\ a $365$-period evaluation window,
matching the protocol whose value-iteration optimum is reproduced cell-for-cell by the in-crate exact
MDP (Section~\ref{sec:perishable-setup}). As for the other problems, $G(\btheta)$ is estimated by
simulation over shared demand-path seeds.

\FloatBarrier

\subsection{Policy}
\label{sec:perishable-policy}

Perishable inventory tests whether the same recipe (Figure~\ref{fig:policy-pipeline}) carries over
when the state is an \emph{age-structured} stock vector and the optimal control is known to be
non-trivially state-dependent (the published optimal policy orders \emph{less} when older stock is
already on hand). \textbf{Input.} The policy input is this problem's age-structured state
$\mathbf S_t=(o^{1}_t,o^{2}_t)$ (youngest and oldest on-hand buckets; for $L>1$ the in-transit
pipeline is prepended), divide-by-scale normalized by the demand mean as in
Section~\ref{sec:methods-principle}. \textbf{Output.} A valid action is a single non-negative integer
order $q_t\in\{0,\dots,\bar q\}$. \textbf{Internals.} Following the policy interface of
Section~\ref{sec:methods-principle}, $\pi_{\btheta}$ applies the soft-tree split/leaf mechanism of
Section~\ref{sec:soft-tree} as its backbone: a depth-2 tree with \emph{oblique} (full-dimensional)
splits and \emph{linear} leaves, so each leaf emits an affine function of the age vector that is
clipped to the order cap. This is the smallest tree that can express the published optimal policy's
``order less when older stock is on hand'' behaviour --- an oblique split separates age profiles and
the linear leaf tilts the order against the on-hand buckets --- while a constant-leaf or single base-
stock level cannot. The whole policy carries $21$ parameters. We warm-start CMA-ES at the encoded
best base-stock level (Section~\ref{sec:methods-cmaes}): the generation-0 mean return reproduces the
base-stock gate to within its standard error, so the search starts \emph{at} the heuristic and any
improvement is attributable to the learned age-dependence rather than to a lucky initialization.

\subsection{Experiment setup}
\label{sec:perishable-setup}

We use the two De Moor et al.\ (2022) Scenario~A instances that the repository can verify exactly:
\texttt{de\_moor2022\_m2\_exp2\_l1\_cp7\_fifo} (FIFO) and its LIFO sibling
\texttt{de\_moor2022\_m2\_exp1\_l1\_cp7\_lifo}, sharing the parameters of
Table~\ref{tab:perish-params}. These are the only two instances small enough ($11^2=121$ states) for
the in-crate exact value-iteration MDP, which re-derives --- not merely stores --- the published
optimal-policy tables, the best base-stock levels ($S=7$ FIFO, $S=5$ LIFO), and the
\citet{farrington2025perishable} Table~3 value-iteration returns ($-1457$ FIFO, $-1553$ LIFO).

\paragraph{Comparator gate.}
The like-for-like comparator is the best single base-stock policy (order-up-to
$q_t=(S-\text{inventory position})^+$) at its exact-optimal level, \emph{re-scored on the same
Monte-Carlo estimator} used to score the learned policy. This is the appropriate ``heuristic to
beat'': beating the best base-stock policy under an identical estimator is a clean apples-to-apples
win, free of the estimator mismatch discussed below.

\paragraph{Evaluation protocol and selection.}
Policies are trained with warm-started CMA-ES and evaluated under the common-random-number protocol of
Section~\ref{sec:methods-eval}. Three \emph{disjoint} seed blocks are used: a search block ($64$
seeds) drives CMA-ES, the best generation is then chosen on a \emph{separate validation block} ($512$
seeds), and the selected policy is reported on a third \emph{held-out evaluation block} ($2048$ seeds,
horizon $465$). Selecting on the disjoint validation block is essential: it yields the gate-beating
policy reported below, whereas selecting the incumbent on the evaluation block itself overfits the
estimator and flips the FIFO result to a small ($\approx-0.49\%$) loss. The base-stock gate is scored
on the identical evaluation block.

\paragraph{The value-iteration optimum is context only.}
The published $-1457$ / $-1553$ optima are the \emph{analytic} expected discounted returns under a
midpoint-binned Gamma demand (matched to Farrington Table~3), whereas the learned and gate rows are
\emph{Monte-Carlo} means over sampled-and-rounded Gamma rollouts. On this slice the optimal base-stock
level evaluates $\approx1\%$ worse under the Monte-Carlo estimator than under the analytic one --- a
systematic offset that is a property of the estimator, not of any policy. We therefore report the
gate beat as the headline result and treat proximity to the VI optimum as informational context.

\begin{table}[!htbp]
\centering
\caption{Perishable benchmark instances of \citet{deMoor2022perishable} used here (the two with an
in-crate exact-MDP verifier). Both have shelf life $m=2$, lead time $L=1$, rounded-Gamma demand with
mean $\mu=4$ and coefficient of variation $0.5$, order cap $\bar q=10$, horizon $465$ with a
$100$-period warm-up, and discount $\gamma=0.99$; they differ only in the issuing rule. Source:
\texttt{src/problems/perishable\_inventory/references.rs}.}
\label{tab:perish-params}
\begin{tabular}{ll}
\toprule
Parameter & Value \\
\midrule
Shelf life $m$ / lead time $L$ & $2$ / $1$ \\
Demand (rounded Gamma) $\mu,\ \mathrm{cv}$ & $4,\ 0.5$ \\
Procurement $c$ / holding $h$ & $3$ / $1$ \\
Lost-sales penalty $p$ / waste $c_w$ & $5$ / $7$ \\
Issuing rule & FIFO \emph{or} LIFO \\
\bottomrule
\end{tabular}
\end{table}
\FloatBarrier

\subsection{Results}
\label{sec:perishable-results}

Table~\ref{tab:perish-results} compares the best learned depth-2 oblique-linear soft tree against the
best base-stock gate under the common-random-number protocol. On both instances the learned policy
\emph{beats} the gate by a margin that is many times its paired standard error: on FIFO it improves
the discounted return from $-1475.08$ to $-1457.90$ ($+1.16\%$, $\approx9.5\times$ the paired SEM),
and on LIFO from $-1565.98$ to $-1553.16$ ($+0.82\%$, $\approx6.6\times$ the paired SEM). Both margins
are statistically clear gate beats, and both were obtained by selecting on the disjoint validation
block (Section~\ref{sec:perishable-setup}).

\paragraph{Interpretation.}
The headline is the \emph{like-for-like gate beat}: under one shared Monte-Carlo estimator, a $21$-
parameter age-dependent soft tree improves on the best base-stock level that the same estimator scores
--- a level that is itself the exact optimum within the base-stock family. The learned tree exploits
the age structure that a single base-stock cannot: it orders less when older stock is already on hand,
trading a little holding for fewer outdated units, which is exactly the published optimal policy's
shape. As context, both learned returns sit essentially \emph{on} the analytic value-iteration optimum
($-1457.28$ FIFO, $-1552.99$ LIFO), but this comparison mixes the analytic and Monte-Carlo estimators
(the gate, scored on the Monte-Carlo estimator, sits $\approx1\%$ below the analytic optimum purely
from that mismatch), so we read the VI proximity as corroborating context rather than as a second,
independent win. The FIFO margin is the larger of the two, consistent with FIFO leaving more room over
the base-stock heuristic than the already near-heuristic-optimal LIFO case.

\begin{table}[!htbp]
\centering
\small
\caption{Perishable inventory: best learned depth-2 oblique-linear soft tree vs.\ the best base-stock
gate, common-random-number evaluation ($2048$ held-out seeds, horizon $465$, $\gamma=0.99$). Returns
are discounted \emph{negated} costs (higher is better). $\Delta\%$ is the learned policy's improvement
over the gate; ``SEM mult.'' is that gain divided by the paired standard error. The value-iteration
optimum column is the \emph{analytic} reference of \citet{deMoor2022perishable} /
\citet{farrington2025perishable} and is context only (different estimator). Source:
\texttt{outputs/autoresearch/perishable\_inventory\_autoresearch/}.}
\label{tab:perish-results}
\begin{tabular}{lcc}
\toprule
Policy / metric & FIFO ($S=7$) & LIFO ($S=5$) \\
\midrule
Best base-stock gate & $-1475.08$ & $-1565.98$ \\
Learned soft tree (ours) & $\mathbf{-1457.90}$ & $\mathbf{-1553.16}$ \\
Improvement over gate $\Delta\%$ & $\mathbf{+1.16\%}$ & $\mathbf{+0.82\%}$ \\
Gain / paired SEM & $\approx 9.5\times$ & $\approx 6.6\times$ \\
\midrule
VI optimum (analytic, context) & $-1457.28$ & $-1552.99$ \\
\bottomrule
\end{tabular}
\par\medskip\footnotesize The gate beat is the like-for-like win: gate and learned policy are scored on
the \emph{same} Monte-Carlo estimator and the gate is the exact optimum within the base-stock family.
The VI-optimum row is the analytic discounted return under midpoint-binned Gamma demand and is not
on the same estimator (the gate sits $\approx1\%$ below it purely from that mismatch); we report it
as context, not as a separate improvement claim. Selection is on a disjoint validation block;
selecting on the evaluation block instead flips FIFO to a small ($\approx-0.49\%$) loss.
\end{table}
\FloatBarrier
